\uuid{C2Sm}
\titre{Noyau/Image/Théorème du rang}
\auteur{hanani ; mzari}
\organisation{exo7}
\licence{CC BY-NC-SA 4.0}
\datecreate{2020-01-20}
\niveau{L1}
\module{Algèbre}
\chapitre{Application linéaire}
\sousChapitre{Définition}
\difficulte{2}
\theme{Noyau et image, Noyau, Image, Théorème du rang}
\exo7qcmid{lille-2.applications-lineaires.018}

\contenu{
\begin{question}
\qclassification{107.01}
\qtags{motcle=Noyau/Image/Théorème du rang}
On considère $\Rr^3$ muni de la base canonique ${\cal{ B}}= \{e_1,e_2,e_3\}$ et $f$ l'endomorphisme de $\Rr^3$ défini par
$f(e_1)= e_3,\, f(e_2)= e_1+e_2,\,  f(e_3)= e_1+e_2+e_3$.
Quelles sont les assertions vraies ?
\begin{answers}
\bad{$\{e_1+e_2-e_3\}$ est une base de $\Im f $}.
\good{$\dim \Im f =2$}.
\good{$\{e_1+e_2-e_3\}$ est une base de $\ker f$}.
\bad{$\dim \ker f=2 $}.
\end{answers}
\begin{explanations}  On remarque que $f(e_3)= f(e_1)+f(e_2) $ et que $f(e_1)$ et $f(e_2)$ ne sont pas colinéaires, donc $\{e_1+e_2,e_3\}$ est une base de $\Im f$ et donc $\dim \Im f =2$.
\vskip0mm
D'après le théorème du rang, $\dim \ker f = 1 $ et comme $f(e_1+e_2-e_3) = 0$ et  $e_1+e_2-e_3 \neq 0$, $\{e_1+e_2-e_3\}$ est une base de $\ker f$.
\end{explanations}
\end{question}
}
